\documentclass[12pt, letterpaper]{article}

%-------Packages---------
\usepackage{amssymb,amsfonts}
\usepackage{mathptmx}
\usepackage{amsthm}
\usepackage{graphicx}
\usepackage[all,arc]{xy}
\usepackage[colorlinks=true, allcolors=blue]{hyperref}
\usepackage{enumerate}
\usepackage{bbm}
\usepackage{dsfont}
\usepackage{mathrsfs}
\usepackage{tikz-cd}
\usepackage{setspace}
\usepackage{import}
\usepackage[utf8]{inputenc}
\usepackage{hyperref}
\usepackage{tikz}
\usepackage{float}
\usepackage{mdframed}
\usepackage[nocheck]{fancyhdr}
\usepackage[nointegrals]{wasysym}

\usepackage[left=1in,top=1in,right=1in,bottom=1in]{geometry}

\usepackage{mathtools}
\usepackage{setspace}
\usepackage{cleveref}
\usepackage{multicol}

%--------Theorem Environments--------
%theoremstyle{plain} --- default
\newmdtheoremenv{theorem}{Theorem}[subsection]
\newmdtheoremenv{corollary}[theorem]{Corollary}
\newtheorem{proposition}[theorem]{Proposition}
\newmdtheoremenv{lemma}[theorem]{Lemma}
\newtheorem{conj}[theorem]{Conjecture}
\newtheorem{quest}[theorem]{Question}

\theoremstyle{definition}
\newmdtheoremenv{definition}[theorem]{Definition}
\newmdtheoremenv{definitions}[theorem]{Definitions}
\newtheorem{example}[theorem]{Example}
\newtheorem{algorithm}[theorem]{Algorithm}
\newtheorem{rmk}[theorem]{Remark}
\newtheorem{exercise}[theorem]{Exercise}

%----------- my commands -------------
\newcommand{\C}{\mathbb{C}}
\newcommand{\R}{\mathbb{R}}
\newcommand{\Q}{\mathbb{Q}}
\newcommand{\Z}{\mathbb{Z}}
\newcommand{\N}{\mathbb{N}}
\renewcommand{\P}{\mathbb{P}}
\newcommand{\contr}{\Rightarrow \Leftarrow}
\newcommand{\HRule}[1]{\rule{\linewidth}{#1}}
\newcommand{\s}{\(\,\)}
\renewcommand{\t}[1]{\text{#1}}
\newcommand{\ang}[1]{\left\langle #1 \right\rangle}

% Meta data
\setstretch{1.2}

\begin{document}

\pagestyle{fancy}
\fancyhf{}
\fancyhead[LOE]{\slshape{\textbf{Vincent Ferrara}}}
\fancyhead[ROE]{\thepage}

\subsection*{Problem 1}
Consider $p_k$ the probability of getting \$20 at $k$ dollars. We have a $0.5$ chance to lose a dollar, $0.4$ 
chance to gain a dollar, and a $0.1$ chance to break even. Thus, $p_k=0.5p_{k-1}+0.4p_{k+1}+0.1p_k$ thus\\
$p_k=5p_{k-1}/9+4p_{k+1}/9$\\
$p_{k+1}=9p_k/4-5p_{k-1}/4$\\
$p_{k+1}-p_k=(5/4)(p_{k}-p_{k-1})$\\
$p_{k+1}-p_k=(5/4)^{k}(p_1-p_0)$\\
$p_{k+1}-p_k=(5/4)^{k}(p_1)$\\
Consider $\sum\limits_{i=0}^k(p_{k}-p_{k-1})$\\
$\sum\limits_{i=1}^k(p_{i}-p_{i-1})=(p_k-p_{k-1})+(p_{k-1}+p_{k-1})+\dots+(p_1-p_0)=p_k-p_0=p_k$ also from above we know $p_i-p_{i-1}=(5/4)^i(p_1)$\\
$p_k=\sum\limits_{i=1}^k(p_{i}-p_{i-1})=(p_1)\sum\limits_{i=1}^k(5/4)^{i-1}=(p_1)\sum\limits_{i=0}^k(5/4)^{i}=(p_1)\dfrac{1-(5/4)^k}{1-(5/4)}$\\
$p_k=(p_1)\dfrac{1-(5/4)^k}{1-(5/4)}$\\
$p_{20}=1=(p_1)\dfrac{1-(5/4)^{20}}{1-(5/4)}$\\
$p_1=\dfrac{1-(5/4)}{1-(5/4)^{20}}$\\
$p_{10}=\dfrac{1-(5/4)}{1-(5/4)^{20}}\dfrac{1-(5/4)^{10}}{1-(5/4)}=\dfrac{1-(5/4)^{10}}{1-(5/4)^{20}}$

\newpage
\subsection*{Problem 2}
$G_1(s)=s/3G_2(s)+s/3G_1(s)+s/3G_0(s)$\\
$G_1(s)=s/3G_2(s)+s/3G_1(s)+s/3$\\ since $G_0(s)=E[s^{X_0}]=E[s^0]=1$\\
$G_1(s)-s/3G_1(s)=s/3(G_2(s)+1)$\\
$G_1(s)(1-s/3)=s/3(G_2(s)+1)$\\
\phantom{text}\\
Given we know that for 1,3 we are both 1 away from the end points and have the same probability of going up, down, or breaking even, we know by symmetry that $X_1=X_3$ thus $G_1(s)=G_3(s)$\\
Thus $G_2(s)=s/3G_3(s)+s/3G_1(s)+s/3G_2(s)$\\
$G_2(s)=2s/3G_1(s)+s/3G_2(s)$\\
$G_2(s)=2s/(3-s)G_1(s)$\\
Thus plugging in $G_2$ we get that\\
$G_3(s)=G_1(s)=\dfrac{3s-s^2}{9-6s-s^2}$\\
$G_2(s)=\dfrac{2s^2}{9-6s-s^2}$

\newpage
\subsection*{Problem 3}
$E[X | \t{starting at } HH]=(1/2)(1+E[X | \t{starting at } HH])+(1/2)$\\
$E[X | \t{starting at } HH]=2$\\
\phantom{te}\\
$E[X]=E[X \mid T]\P(T)+E[X \mid HT]\P(HT)+E[X \mid HHT]\P(HHT)+E[X \mid HHH]\P(HHH)$\\
$E[X]=1/2E[X+1]+1/4E[X+2]+E[3]1/8+(3+E[X | \t{starting at } HH])1/8$\\
$E[X]=1/2E[X]+1/2+1/4E[X]+1/2+3/8+3/8+2/8$\\
$1/4E[X]=2$\\
$E[X]=8$

\newpage
\subsection*{Problem 4}
$G_n(s)=spG_{n-1}+s(1-p)G_{n}$\\
$G_n(s)=\left( \dfrac{sp}{1-(p-1)s} \right)G_{n-1}$\\
$G_n(s)=\left( \dfrac{sp}{1-(p-1)s} \right)^nG_{0}$\\
$G_0(s)=1$ since there is one way to flip zero heads\\
$G_n(s)=\left( \dfrac{sp}{1-(p-1)s} \right)^n$\\
$G_n(s)=s^np^n\left( \dfrac{1}{1-(p-1)s} \right)^n$\\
$G_n(s)=s^np^n\sum\limits_{k=0}^\infty(-1)^k\binom{n+k-1}{k}(-(1-p)s)^k$\\
$G_n(s)=s^np^n\sum\limits_{k=0}^\infty\binom{n+k-1}{k}(1-p)^ks^k$\\
$G_n(s)=\sum\limits_{k=0}^\infty\binom{n+k-1}{k}p^n(1-p)^ks^{k+n}$\\
Let $k=k+n$\\
$G_n(s)=\sum\limits_{k=n}^\infty\binom{k-1}{n-1}p^n(1-p)^{k-n}s^{k}$\\
Thus $\P(X_n=k)=\binom{k-1}{n-1}p^n(1-p)^{k-n}$ for $k \geq n$ and $0$ for $k < n$
\end{document}
